\documentclass[11pt]{article}

% --- arXiv-friendly preamble (standard packages only) ---
\usepackage[utf8]{inputenc}
\usepackage[T1]{fontenc}
\usepackage{amsmath,amssymb,amsthm}
\usepackage{booktabs}
\usepackage{geometry}
\usepackage{hyperref}
\usepackage{xcolor}
\geometry{margin=1in}
\hypersetup{colorlinks=true,linkcolor=blue,citecolor=blue,urlcolor=blue}

\newtheorem{definition}{Definition}
\newtheorem{proposition}{Proposition}
\newtheorem{theorem}{Theorem}

\newcommand{\dcausal}{d_{\mathrm{c}}}
\newcommand{\Nk}{\mathcal{N}_k}
\newcommand{\code}[1]{\texttt{#1}}

\title{\textbf{Causal Context Regions:\\ A Spatial Memory Model for Long-Horizon LLM Coding Agents}}
\author{The CCOS Project\thanks{Causal Context Operating System (CCOS), an open
research prototype. Source, reproduction scripts and this paper:
\url{https://github.com/CHECKUPAUTO/CCOS}.}}
\date{\today}

\begin{document}
\maketitle

\begin{abstract}
Coding agents backed by large language models (LLMs) drown in context: a
repository is paged into a bounded window as a flat, ranked list of files, and as
tasks lengthen the window fills with globally ``important'' but locally
irrelevant material. We argue that the unit of context should not be the file but
the \emph{region}: a causally coherent neighbourhood of code, dependencies and
associated faults that is hydrated as a whole, the way an operating system pages
in a working set rather than a single address. We present the \emph{Context
Region Engine} (CCOS v0.3), which lifts a 1-D causal memory graph into a spatial
map of regions with a temperature, a causal density and a dynamic admission
policy. We give a precise, deterministic definition of a region in terms of
connected components of a cross-file causal-link graph, and of causal distance as
a weighted shortest path. We prove that regions are a pure function of the graph
and therefore reconstruct bit-for-bit on replay, preserving CCOS's auditability
invariant. On the system's own source tree we measure, \emph{without any LLM
calls}, that region-based selection covers a task's $k$-hop causal neighbourhood
with recall $0.97$ versus $0.35$ for the flat top-score strategy, at
$\approx\!48\%$ fewer tokens for equal coverage, with regions that are
$95.5\%$ internally connected. In a deterministic, LLM-free \emph{simulation}
under an explicit retrieval oracle, classical (lexical) RAG solves $0\%$ of
cross-file causal tasks while structure-aware regional selection solves $100\%$,
with the gap widening with task diameter --- though a strong graph-traversal
baseline ties it on a clean query. Under a \emph{misleading} query, however, every
lexically-seeded method --- including that graph baseline and an ablation of CCOS
that trusts the query --- collapses to $0\%$, while only the workspace-anchored
region survives: the differentiator is the \emph{anchor} (a structural workspace
signal vs.\ a lexical query), not the region machinery. This tests the
\emph{necessary} (retrieval) condition, not the \emph{sufficient} (generation)
one; a first on-device real-LLM run (\code{qwen2.5-7B} via Ollama on an NVIDIA
Jetson AGX Thor) then reproduces the coverage result on real file content and
exhibits the same noisy separation --- under a misleading query only the
workspace-anchored region yields any task success --- while we are explicit that
the small model's arithmetic reasoning, not selection, caps absolute success. We
further specify a falsifiable protocol against classical RAG, GraphRAG, MemGPT and
LangGraph on long-horizon tasks (SWE-bench and a multi-file bug suite). We label
throughout what is \emph{proven} (formal structure, determinism, locality),
\emph{measured} (retrieval coverage, in simulation and on a real LLM), and
\emph{still hypothesised} (large-model and external-benchmark agent gains), so the
contribution can be evaluated honestly.
\end{abstract}

\section{Introduction}

An LLM coding agent operating over a repository must repeatedly decide \emph{what
to put in the context window}. The dominant pattern is retrieval: embed code
chunks, rank them against the task, and concatenate the top-$k$ until a token
budget is exhausted~\cite{lewis2020rag}. This treats context as a one-dimensional
ranked list. Two failure modes follow. First, \emph{dilution}: as tasks lengthen,
the window accumulates globally salient but task-irrelevant code, and models
attend poorly to the middle of long contexts~\cite{liu2023lost}. Second,
\emph{incoherence}: the top-$k$ chunks need not form a connected unit of
reasoning --- a function may be paged in without its callers, its error sites or
the data it depends on.

Operating systems faced the analogous problem decades ago and answered with the
\emph{working set}: the set of recently referenced pages that must stay
resident~\cite{denning1968working}. CCOS adopts the analogy for LLM
context~\cite{packer2023memgpt}: source is parsed into a causal
graph, nodes are scored, and a bounded ``context window'' is paged in and out
like RAM~$\leftrightarrow$~VRAM. CCOS v0.2 keeps this at one dimension: a scalar
score per node. The present work (v0.3) makes the memory \emph{spatial}. Our
contributions are:

\begin{enumerate}
  \item A \textbf{formal model of context regions} (\S\ref{sec:regions}):
  causal distance as a weighted shortest path, and region membership as a
  connected component of a cross-file causal-link graph, with external
  dependency hubs collapsed.
  \item A \textbf{dynamic admission policy} (\S\ref{sec:policy}) that replaces a
  global threshold with a per-region score modulated by temperature, causal
  density, token pressure and task complexity.
  \item A \textbf{determinism theorem} (\S\ref{sec:determinism}): regions are a
  pure function of the graph, and a session reconstructs bit-for-bit from its
  event log, preserving tamper-evident replay.
  \item An \textbf{LLM-free locality evaluation} (\S\ref{sec:eval}) with real,
  reproducible numbers, plus a \textbf{falsifiable comparison protocol}
  (\S\ref{sec:protocol}) against RAG, GraphRAG, MemGPT and LangGraph that we
  state but do not yet run.
\end{enumerate}

\section{Related Work}\label{sec:related}

\paragraph{Retrieval-augmented generation.} RAG augments a parametric model with
a non-parametric store and retrieves passages per
query~\cite{lewis2020rag}. Self-RAG adds reflection tokens that gate
retrieval and critique generations~\cite{asai2023selfrag}. These operate on
\emph{independent} chunks; coherence between retrieved items is not modelled.

\paragraph{Graph- and structure-aware retrieval.} GraphRAG builds an entity
knowledge graph and answers global queries by summarising community
structure~\cite{edge2024graphrag}. For code, property graphs unify AST, control
flow and data flow~\cite{yamaguchi2014cpg}. CCOS regions are in this lineage but
target \emph{paging}: which connected sub-structure to make resident for a task,
under a token budget, with deterministic eviction.

\paragraph{Agent memory.} MemGPT casts the LLM as an OS that pages between an
in-context window and external storage~\cite{packer2023memgpt}; Generative Agents
retrieve memories scored by recency, importance and
relevance~\cite{park2023generative} --- the same factors CCOS aggregates into a
region temperature. LangGraph~\cite{langgraph} structures agents as stateful
graphs of steps; it orchestrates \emph{control} flow, whereas CCOS structures
\emph{memory}. The two are complementary.

\paragraph{Context-window management.} KV-cache eviction keeps ``heavy hitter''
or sink tokens resident~\cite{liu2023lost}, and PagedAttention applies OS paging
to the KV cache~\cite{kwon2023paged}. These act at the token level inside one
forward pass; CCOS acts at the \emph{semantic} level across an agent session.

\section{The Causal Context Substrate}\label{sec:substrate}

CCOS parses source into a directed \emph{causal memory graph} $G=(V,E)$. A node
$v\in V$ is a file, module, symbol or external dependency, with scalar fields
$\mathrm{imp}(v)$ (base importance), $\mathrm{fail}(v)\in[0,1]$ (failure
relevance) and $\mathrm{rec}(v)\in[0,1]$ (recency). An edge $e=(u\!\to\!w)\in E$
carries a weight $w(e)\in(0,1]$ and a type (containment, dependency, reference,
causation). The kernel assigns each node a causal score
\begin{equation}
\mathrm{score}(v) = \mathrm{clamp}\big(0.15\,\mathrm{imp}(v) + 0.50\,\mathrm{fail}(v)
  + 0.30\,\mathrm{rec}(v) + 0.05\ln(1{+}\mathrm{acc}(v)),\,0,1\big),
\label{eq:score}
\end{equation}
where $\mathrm{acc}(v)$ is the access count. Faults propagate along edges,
$\mathrm{fail}$ decays with a logical clock, and every state transition is
appended to a hash-chained event log enabling deterministic
replay~\cite{haber1991timestamp}. Node identifiers are namespaced
(\code{file:p}, \code{mod:p:n}, \code{use:p:path}, \code{sym:p:n},
\code{dep:root}); the owning file of a node is recoverable from its identifier.
CCOS v0.2 pages in the top-$\mathrm{score}$ nodes: a flat, 1-D policy. We now
make selection spatial.

\section{Context Regions}\label{sec:regions}

\subsection{Causal distance}

\begin{definition}[Causal distance]
Let $\hat G$ be the undirected multigraph on $V$ induced by $E$, and assign each
edge the cost $c(e) = -\ln w(e) \ge 0$, so that a stronger causal link is a
shorter step. The \emph{causal distance} $\dcausal(u,v)$ is the minimum total
cost over all $u$--$v$ paths in $\hat G$, and $+\infty$ if none exists. The
unweighted hop distance $\mathrm{hops}(u,v)$ is defined analogously with unit
costs.
\end{definition}

$c(e)\ge 0$ since $w(e)\le 1$, so $\dcausal$ is a genuine shortest-path metric on
each connected component (non-negative, symmetric, triangle inequality). The
\emph{$k$-hop causal neighbourhood} of a target $t$ is
\begin{equation}
\Nk(t) = \{\, v \in V : \mathrm{hops}(t,v) \le k \,\}.
\end{equation}
$\Nk(t)$ is the ground truth ``what is causally relevant to a task at $t$'' used
in the evaluation (\S\ref{sec:eval}).

\subsection{Region membership}

External dependency hubs (e.g.\ \code{dep:std}) connect almost everything and
must not collapse the graph into one component. We therefore separate them.

\begin{definition}[Cross-file causal link]
For a non-external node $v$ let $\phi(v)$ be its owning file. Two files $f,g$ are
\emph{directly linked}, written $f \approx g$, iff there exists an edge
$(u\!\to\!w)\in E$ with $\phi(u)=f$, $\phi(w)=g$, $f\neq g$, and neither $u$ nor
$w$ is an external dependency node.
\end{definition}

\begin{definition}[Region]\label{def:region}
Let $\approx^{*}$ be the reflexive-transitive closure of $\approx$ on the set of
file keys. A \emph{region} is the set of all nodes whose files lie in one
equivalence class of $\approx^{*}$; all external dependency nodes form one
additional region. Two nodes belong to the same region iff their files are in the
same connected component of the cross-file causal-link graph.
\end{definition}

\begin{proposition}[Regions partition the nodes]
The regions of Definition~\ref{def:region} form a partition of $V$: every node
lies in exactly one region.
\end{proposition}
\begin{proof}
$\approx^{*}$ is an equivalence relation on file keys, so its classes partition
the file keys; the external bucket is a distinct class by construction. Mapping
each non-external node to its file's class and each external node to the external
class is a total function into disjoint blocks, hence a partition. \qed
\end{proof}

By default (only containment and import edges) each source file is its own region.
A genuine cross-file dependency or a propagated failure merges files into a single
multi-file region --- the ``zone of knowledge'' an agent should wake together.

\subsection{Region scalars}

For a region $R$ with member set $M$:
\begin{align}
\mathrm{heat}(v)        &= 0.5\,\mathrm{score}(v) + 0.3\,\mathrm{fail}(v) + 0.2\,\mathrm{rec}(v), \\
\mathrm{temp}(R)        &= \mathrm{clamp}\!\Big(\tfrac{1}{|M|}\textstyle\sum_{v\in M}\mathrm{heat}(v),\,0,1\Big), \label{eq:temp}\\
\mathrm{dens}(R)        &= \frac{|\{\,e\in E : \text{both endpoints} \in M\,\}|}{|M|}. \label{eq:dens}
\end{align}
Temperature is how ``awake'' a region is; density is its internal causal cohesion
(internal edges per member). Activation warms a region
($\mathrm{temp}\mathrel{+}=0.25$, capped) and records a logical tick; a cooldown
step multiplies temperatures by a decay factor and evicts regions below a floor.

\subsection{Dynamic admission policy}\label{sec:policy}

The static threshold $0.6$ becomes a function of token pressure $u\in[0,1]$ (the
used fraction of the budget) and task complexity $\kappa\in[0,1]$:
\begin{align}
\theta(u,\kappa)       &= \mathrm{clamp}(0.6 + 0.3\,u - 0.2\,\kappa,\; 0.05,\; 0.95), \\
a(R)                   &= 0.55\,\mathrm{temp}(R) + 0.30\,\frac{\mathrm{dens}(R)}{1+\mathrm{dens}(R)} + 0.15\,\kappa, \\
\mathrm{admit}(R)      &\iff a(R) \ge \theta(u,\kappa).
\end{align}
A hot, cohesive region can be admitted even when the static $0.6$ would reject it;
a nearly-full window raises $\theta$ so only the hottest regions enter.

\section{Determinism and Replay}\label{sec:determinism}

\begin{theorem}[Regional determinism]
The region partition, every region scalar in
\eqref{eq:temp}--\eqref{eq:dens}, and the activation history are pure,
order-independent functions of the graph $G$ and the sequence of (logically
timestamped) activation events. Consequently, given the event log of a session,
the engine state reconstructs identically: if $G'$ is the graph rebuilt from the
log and $L$ its region events, then $\mathrm{replay\_from}(G',L)$ equals the live
engine that produced $L$.
\end{theorem}
\begin{proof}[Proof sketch]
Clustering enumerates nodes and edges in sorted order and computes connected
components by a sorted breadth-first search, so its output is independent of hash
iteration order. Equations \eqref{eq:score}, \eqref{eq:temp} and \eqref{eq:dens}
are deterministic arithmetic over node fields and a count of qualifying edges.
Activation reads a logical clock (a counter), never wall-clock time, and the
emitted \code{RegionActivated} event records the exact tick. Reconstruction
re-clusters $G'$ (identical base state, since $G'\!=\!G$ structurally) and applies
the recorded activations in log order; each step is the same pure function of the
same inputs, so the resulting state is identical. The integration test
\code{replay\_reconstructs\_identical\_engine} checks
$\mathrm{engine}=\mathrm{replay\_from}(G',L)$ by structural equality. \qed
\end{proof}

This extends CCOS's existing guarantees: the primary event log is hash-chained and
tamper-evident, so the region history is auditable, and a 10\,000-cycle test
exhibits no drift in region count or temperatures.

\section{Implementation}\label{sec:impl}

The engine is $\approx\!600$ lines of dependency-light Rust layered above the
graph, with no change to the parser, guard, incremental builder or event-sourcing
core. The public surface is: \code{ContextRegionEngine} (\code{cluster\_nodes},
\code{initialize\_regions}, \code{activate\_region}, \code{tick\_cooldown},
\code{replay\_from}), the data types \code{ContextRegion} / \code{ContextPoint} /
\code{ContextWindow}, the \code{ContextPolicy}, and five new event variants
(\code{RegionCreated}, \code{RegionActivated}, \code{RegionMerged},
\code{RegionEvicted}, \code{ContextWindowGenerated}). A \code{ccos regions} CLI
exposes clustering, activation and the locality report. The whole crate builds
warning-clean under \code{clippy -D warnings} and passes 202 tests.

\section{Evaluation: what we can measure today}\label{sec:eval}

We separate \emph{measured} results (this section, fully reproducible and
LLM-free) from \emph{hypothesised} agent-level gains (\S\ref{sec:protocol}).

\paragraph{Setup.} We run on CCOS's own source tree
($|V|=705$ nodes, $|E|=822$ edges, yielding $26$ regions). For a target node $t$
we take $\Nk(t)$ with $k=2$ as ground truth and compare two selection strategies
at the region's size budget: \textbf{flat} --- the top-$|R|$ nodes by
$\mathrm{score}$ (CCOS v0.2 / classical ranked retrieval) --- and \textbf{region}
--- the members of $t$'s region. We report causal precision $|S\cap\Nk|/|S|$,
recall $|S\cap\Nk|/|\Nk|$, and the tokens each strategy needs to \emph{cover}
$\Nk$. All numbers are emitted by \code{scripts/region\_benchmark.sh} and are
deterministic.

\begin{table}[h]
\centering
\begin{tabular}{lcc}
\toprule
Strategy & causal precision & causal recall \\
\midrule
flat (v0.2 / ranked) & 0.021 & 0.347 \\
\textbf{region (v0.3)} & \textbf{0.057} & \textbf{0.972} \\
\bottomrule
\end{tabular}
\caption{Locality on \code{src/} ($k=2$, 12 targets). Region selection covers
$97\%$ of a task's causal neighbourhood versus $35\%$ for flat, at
$\approx\!48\%$ fewer tokens for equal coverage.}
\label{tab:locality}
\end{table}

\paragraph{Cohesion and cost.} Regions reach an \textbf{average causal density of
$0.955$} (Eq.~\ref{eq:dens}, normalised) --- they are almost entirely internally
connected, i.e.\ genuine causal clusters rather than arbitrary file groups.
Building the region map for $705$ nodes takes $\approx\!20$\,ms
($54.5$ builds/s); a single activation is $\approx\!20$\,ms.

\paragraph{Honest reading.} Absolute precision is low ($0.06$) because the v0.2
parser emits only containment and import edges, so $\Nk(t)$ is tiny (often $2$--$3$
nodes) while a region spans a whole file. The robust, real wins are
\emph{recall} ($0.97$ vs $0.35$) and \emph{token efficiency} ($\approx\!48\%$):
a region reliably contains a task's causal neighbourhood and pays fewer tokens to
do so. Richer semantic edges (call graph / data flow, roadmap item P1.3) would
sharpen $\Nk$ and is the main lever to raise precision; we flag this rather than
hide it.

\section{Hypothesis simulation under a stated oracle (LLM-free)}\label{sec:sim}

The core thesis --- \emph{does regional causal memory help an agent on long,
multi-file tasks?} --- ultimately needs LLM rollouts (\S\ref{sec:protocol}). But
its \emph{necessary condition} is testable now, without an LLM: \textbf{an agent
cannot solve a task whose required causal context is absent from its window}. We
measure, under an explicit oracle, whether each selection strategy \emph{places
the required causal context in the window}.

\paragraph{Setup.} We generate modular synthetic repositories: $S$ independent
subsystems, each a set of files linked by one cross-file causal chain plus
high-score \emph{decoy} symbols; there are no edges between subsystems, so each
subsystem is one bounded causal region. The causal structure is \emph{decoupled}
from the lexical structure --- a chain neighbour shares no identifier token with
the target, only an edge --- modelling a dependency that is causally essential yet
lexically dissimilar. A task of \emph{diameter} $d$ requires the $\pm d$ window
along its chain (up to $2d{+}1$ files); the budget holds $\approx$ one subsystem,
far less than the repository. The \textbf{oracle} is $R(t)\subseteq S$.

Crucially, retrieval pipelines locate code from a text \emph{query}, whereas an
OS-style memory \emph{anchors} on the workspace signal (the active file, a failing
test). We model both: each task carries a query (a token bag) and an anchor (the
active file node), and we run two scenarios --- \textbf{clean} (the query points
at the target) and \textbf{noisy} (a trap decoy in an unrelated subsystem
out-scores the target lexically). Six strategies share the budget:
\code{rag-dense} (top-$k$ lexical), \code{rag-hybrid} (lexical $+$ causal score),
\code{graphrag-1hop} (best hit $+$ one hop), \code{graphrag-bfs} (unbounded
expansion from the best hit), \code{ccos-from-query} (CCOS region of the best
\emph{lexical} hit --- an ablation), and \code{ccos-region} (CCOS region of the
\emph{anchor}). Seeded and deterministic; reproduce with \code{ccos experiment}.

\begin{table}[h]
\centering
\begin{tabular}{lcccccc}
\toprule
& \multicolumn{4}{c}{success at diameter $d$} & \multicolumn{2}{c}{overall} \\
\cmidrule(lr){2-5}\cmidrule(lr){6-7}
Strategy & $d{=}1$ & $d{=}2$ & $d{=}3$ & $d{=}4$ & succ. & cov. \\
\midrule
\multicolumn{7}{l}{\emph{Clean query (points at the target)}}\\
\code{rag-dense}      & 0.00 & 0.00 & 0.00 & 0.00 & 0.00 & 0.19 \\
\code{rag-hybrid}     & 1.00 & 0.00 & 0.00 & 0.00 & 0.23 & 0.65 \\
\code{graphrag-1hop}  & 1.00 & 0.00 & 0.00 & 0.00 & 0.23 & 0.58 \\
\code{graphrag-bfs}   & 1.00 & 1.00 & 1.00 & 1.00 & 1.00 & 1.00 \\
\code{ccos-from-query}& 1.00 & 1.00 & 1.00 & 1.00 & 1.00 & 1.00 \\
\code{ccos-region}    & 1.00 & 1.00 & 1.00 & 1.00 & 1.00 & 1.00 \\
\midrule
\multicolumn{7}{l}{\emph{Noisy query (a decoy out-scores the target lexically)}}\\
\code{rag-dense}      & 0.00 & 0.00 & 0.00 & 0.00 & 0.00 & 0.19 \\
\code{rag-hybrid}     & 1.00 & 0.00 & 0.00 & 0.00 & 0.23 & 0.65 \\
\code{graphrag-1hop}  & 0.00 & 0.00 & 0.00 & 0.00 & 0.00 & 0.19 \\
\code{graphrag-bfs}   & 0.00 & 0.00 & 0.00 & 0.00 & 0.00 & 0.00 \\
\code{ccos-from-query}& 0.00 & 0.00 & 0.00 & 0.00 & 0.00 & 0.00 \\
\textbf{\code{ccos-region}} & \textbf{1.00} & \textbf{1.00} & \textbf{1.00} & \textbf{1.00} & \textbf{1.00} & \textbf{1.00} \\
\bottomrule
\end{tabular}
\caption{Hypothesis simulation ($800$ tasks, seed $42$, budget $\approx$ one
subsystem). Success $=$ the required causal set $R(t)$ is inside the window. Under
a clean query every structure-aware method ties; under a misleading query only
the workspace-anchored region survives.}
\label{tab:sim}
\end{table}

\paragraph{Findings (and their honest limits).}
\textbf{(1)} Lexical retrieval (\code{rag-dense}) \emph{fails} on cross-file
causal tasks ($0\%$; coverage $0.19$): similarity cannot surface
causally-essential but lexically-dissimilar context --- the hypothesis's premise.
(\code{rag-hybrid}'s $d{=}1$ wins come from the \emph{causal score} it borrows,
not lexical similarity, which is why they persist under noise.)
\textbf{(2)} The value of structure-aware selection \emph{grows with diameter}:
one-hop expansion solves only $d{=}1$, while full causal paging
(\code{graphrag-bfs}, \code{ccos-region}) solves all $d$ --- the direction of H2.
\textbf{(3) The clean tie.} Under a clean query, \code{graphrag-bfs},
\code{ccos-from-query} and \code{ccos-region} all reach $1.00$: the lever is
causal \emph{structure}, \textbf{not CCOS per se}.
\textbf{(4) The noisy separation.} Under a \emph{misleading} query, every method
that locates code lexically collapses to $0\%$ --- including the strong
\code{graphrag-bfs} \emph{and} the \code{ccos-from-query} ablation (CCOS seeded on
the query). Only \code{ccos-region}, which anchors on the workspace signal rather
than the query, survives at $1.00$. The ablation isolates the differentiator: it
is the \emph{anchor source} (a structural workspace signal vs.\ a lexical query),
not the region machinery. This is the realistic regime --- a task description
rarely names the exact distant cause, and an OS-style memory that tracks the
active working set is robust where query-time retrieval is misled.
\textbf{(5) Assumptions, stated.} This credits CCOS with a reliable anchor (the
active file / failing test), which an OS-level memory has but a pure retrieval
pipeline does not; and it assumes \emph{modular} repositories with separable
regions (density $0.955$ on real code, \S\ref{sec:eval}) --- a monolithic giant
region exceeding the budget collapses the advantage (observed, reported).
\textbf{(6)} Throughout this is a \emph{simulation under a stated oracle}: it
tests the \emph{necessary} (retrieval) condition, not the \emph{sufficient}
(generation) one of \S\ref{sec:protocol}.

\section{Real-LLM evaluation: a first on-device measurement}\label{sec:protocol}

The simulation (\S\ref{sec:sim}) established the necessary condition. The
sufficient condition --- that an LLM agent then \emph{solves} more tasks with
regional memory than with chunk retrieval --- requires model rollouts. We
\textbf{implement the harness} (\code{ccos eval}) and report a \textbf{first
real-LLM run} against a local model.

\paragraph{Implemented harness.} Each task is a tiny multi-file project encoding
an \emph{arithmetic causal chain}: a base constant in one file is transformed
through a chain of one-line functions across separate files, and the question
``what integer does the last function return?'' is answerable \emph{only} by
reading the whole chain (the distant cause included). Grading is therefore
exact-match on an integer --- objective and automatable, with no code execution.
The same six strategies of \S\ref{sec:sim} assemble the context window under a
token budget (with the clean/noisy query split), the window is sent to a model
(any OpenAI-, Anthropic-Messages-, or Ollama-compatible endpoint), and we record
three metrics per diameter: \textbf{task success} (correct integer),
\textbf{oracle coverage} (required chain $\subseteq$ window --- model-independent),
and \textbf{symbol-hallucination} (the answer cites a function absent from the
project).

\paragraph{Setup.} We run \code{ccos eval} on-device on an NVIDIA Jetson AGX Thor
against \code{qwen2.5:7b-instruct} served by Ollama (temperature $0$), with $20$
tasks, seed $7$, a $600$-token budget, diameters $1$--$4$, in both the clean and
noisy regimes. Table~\ref{tab:real} reports the run.

\begin{table}[h]
\centering
\begin{tabular}{lcccccr}
\toprule
& \multicolumn{4}{c}{success at diameter $d$} & & \\
\cmidrule(lr){2-5}
Strategy & $d{=}1$ & $d{=}2$ & $d{=}3$ & $d{=}4$ & cov. & tok. \\
\midrule
\multicolumn{7}{l}{\emph{Clean query (names the target function)}}\\
\code{rag-dense}      & 0.12 & 0.00 & 0.00 & 0.00 & 0.00 & 519 \\
\code{rag-hybrid}     & 0.00 & 0.00 & 0.00 & 0.00 & 0.00 & 521 \\
\code{graphrag-1hop}  & 1.00 & 0.00 & 0.00 & 0.00 & 0.40 & 270 \\
\code{graphrag-bfs}   & 1.00 & 0.00 & 0.17 & 0.00 & 1.00 & 402 \\
\code{ccos-from-query}& 1.00 & 0.00 & 0.17 & 0.00 & 1.00 & 402 \\
\code{ccos-region}    & 1.00 & 0.00 & 0.17 & 0.00 & 1.00 & 402 \\
\midrule
\multicolumn{7}{l}{\emph{Noisy query (a decoy out-scores the target lexically)}}\\
\code{rag-dense}      & 0.00 & 0.00 & 0.00 & 0.00 & 0.00 & 531 \\
\code{rag-hybrid}     & 0.00 & 0.00 & 0.00 & 0.00 & 0.00 & 521 \\
\code{graphrag-1hop}  & 0.00 & 0.00 & 0.00 & 0.00 & 0.00 & 165 \\
\code{graphrag-bfs}   & 0.00 & 0.00 & 0.00 & 0.00 & 0.00 & 165 \\
\code{ccos-from-query}& 0.00 & 0.00 & 0.00 & 0.00 & 0.00 & 165 \\
\textbf{\code{ccos-region}} & \textbf{1.00} & \textbf{0.00} & \textbf{0.17} & \textbf{0.00} & \textbf{1.00} & \textbf{402} \\
\bottomrule
\end{tabular}
\caption{First real-LLM run: \code{qwen2.5:7b-instruct} via Ollama, on-device
(NVIDIA Jetson AGX Thor), $20$ tasks, seed $7$, budget $600$ tokens. ``cov.'' is
the model-independent oracle coverage; ``tok.'' the mean input tokens. Success is
an exact-match integer answer.}
\label{tab:real}
\end{table}

\paragraph{Findings.}
\textbf{(1) Coverage transfers from simulation to real text.} The
model-independent coverage column reproduces Table~\ref{tab:sim} on real file
content: lexical RAG covers $0$, the structure-aware methods cover $1.00$ under a
clean query, and under noise \emph{only} \code{ccos-region} holds $1.00$ --- the
necessary condition, confirmed outside the simulator.
\textbf{(2) Success is bounded by the model, not the context.} Where coverage is
$1.00$, the $7$B model still solves only $d{=}1$ ($1.00$) and part of $d{=}3$
($0.17$), missing $d{=}2,4$: with the whole chain in the window, the residual
failure is \emph{arithmetic reasoning} --- a model-capability ceiling, not a
selection failure. Crucially no method ever succeeds \emph{without} coverage
(success $\subseteq$ coverage), which is the claim the harness is built to test.
\textbf{(3) The noisy separation holds on a real LLM.} Under a misleading query
every query-driven method collapses to $0\%$ success --- including
\code{graphrag-bfs} and the \code{ccos-from-query} ablation --- while
\code{ccos-region}, anchored on the workspace signal, is the \emph{only} strategy
with non-zero success ($d{=}1$ at $1.00$). The query-driven methods even look
\emph{cheaper} under noise ($165$ vs $402$ tokens) precisely because they
confidently retrieve the wrong, smaller file; \code{ccos-region} spends its budget
on the causally-correct chain.
\textbf{(4) Honest limits.} $20$ tasks ($\approx 5$ per diameter) is a small
sample, and a $7$B model floors the sufficient condition. A frontier model
(\code{deepseek-v4-pro}, in progress) or a $70$B local model is expected to lift
success wherever coverage is $1.00$, \emph{sharpening} --- not changing --- the
separation, since coverage already fixes the ceiling.

\paragraph{Toward external benchmarks.} The arithmetic-chain suite isolates the
selection question; the decisive external tests are SWE-bench~\cite{jimenez2023swebench}
issue resolution (a patch passes the hidden tests) and a controlled
\emph{multi-file bug} suite where the fault site, its cause and its blast radius
lie in distinct files. Baselines share the base LLM and token budget: classical
RAG~\cite{lewis2020rag} (top-$k$ cosine over chunks), GraphRAG~\cite{edge2024graphrag}
(community summaries over a code graph), MemGPT~\cite{packer2023memgpt} (OS-style
paged memory), a LangGraph~\cite{langgraph} agent with a vector store, and CCOS
regions. Metrics: resolved-rate, input tokens to success, and a citation-grounding
hallucination rate checked against the ground-truth graph. We have begun this on
real history: \code{scripts/causal\_validation} mines fix commits, checks out the
pre-fix tree, injects the fault, and scores $R_{\mathrm{cov}} = |F_{\mathrm{target}}
\cap \mathrm{WorkingSet}_K| / |F_{\mathrm{target}}|$ --- the fraction of the files
a fix touched that the bounded working set recovers. On this repository's history
the first run exposed a limitation and its fix, both measured: with
downstream-only failure propagation $R_{\mathrm{cov}}$ was flat at $0.33$ (only the
seed file recovered, as co-changed files are \emph{upstream} importers reached
only through dependency hubs). Resolving intra-crate imports into file$\to$file
edges and propagating failure \emph{bidirectionally} fixes this. On a mature
external codebase --- \code{sharkdp/fd}, $n{=}25$ mined fix commits --- the effect
is unmistakable: $R_{\mathrm{cov}}$ rises from $0.50$ (downstream-only) to $0.92$
($K{=}50$) and $0.96$ ($K{=}100$, geometric mean $0.95$, $92\%$ of bugs fully
covered), while \emph{diluting} to $0.28$ at a tight $K{=}20$ budget --- a real,
falsifiable trade-off. So on real bugs the \emph{necessary} condition (the files a
fix must touch lie in the window) now holds for the large majority; running the
protocol across several repositories, and end-to-end with the LLM, is future work.

\paragraph{Hypotheses.}
\textbf{H1} (efficiency): CCOS reaches equal task success at fewer input tokens
than RAG/GraphRAG, because a region covers the causal neighbourhood with higher
recall (Table~\ref{tab:locality}).
\textbf{H2} (long-horizon success): on multi-file tasks CCOS's resolved-rate
exceeds chunk RAG by a margin that grows with the causal diameter of the task.
\textbf{H3} (grounding): CCOS lowers the symbol-hallucination rate, because the
admitted window is a connected sub-graph of \emph{real} nodes.

\paragraph{Threats to validity.} Region quality is bounded by edge quality
(\S\ref{sec:eval}); results would conflate the region \emph{policy} with the
underlying \emph{graph}. Ablations must therefore fix the graph and vary only the
selection strategy, and report per-task causal diameter so that gains are
attributed to coherence, not to retrieving more text. The \code{ccos-from-query}
ablation in Table~\ref{tab:real} is exactly this control --- same graph and
budget, only the anchor swapped --- and it collapses with the lexical baselines
under noise. A positive result on H1--H3 would constitute the research
contribution; a null result would still validate the deterministic, auditable
infrastructure.

\section{Limitations}

The parser is a line/AST heuristic, so causal neighbourhoods are shallow; regions
are file-granular by default and only merge on explicit cross-file edges;
temperature/policy constants are hand-tuned, not learned; and the agent-level
claims of \S\ref{sec:protocol} are untested. None of these affect the proven
properties (partition, determinism, replay) or the measured locality.

\section{Conclusion}

We reframed LLM context selection from a 1-D ranked list to a spatial map of
\emph{causal regions}, gave the construction a precise and deterministic
definition, proved it reconstructs bit-for-bit on replay, and measured --- without
any LLM --- that regions cover a task's causal neighbourhood far more completely
and cheaply than flat ranking. Whether this locality translates into better
long-horizon agent behaviour is an empirical question we have framed as a
falsifiable protocol against RAG, GraphRAG, MemGPT and LangGraph. We release the
engine, the metrics and this paper so the claim can be tested rather than
asserted.

\begin{thebibliography}{99}
\bibitem{lewis2020rag} P.~Lewis et al. Retrieval-Augmented Generation for
Knowledge-Intensive NLP Tasks. \emph{NeurIPS}, 2020. arXiv:2005.11401.
\bibitem{asai2023selfrag} A.~Asai et al. Self-RAG: Learning to Retrieve, Generate
and Critique through Self-Reflection. \emph{ICLR}, 2024. arXiv:2310.11511.
\bibitem{edge2024graphrag} D.~Edge et al. From Local to Global: A Graph RAG
Approach to Query-Focused Summarization. 2024. arXiv:2404.16130.
\bibitem{yamaguchi2014cpg} F.~Yamaguchi et al. Modeling and Discovering
Vulnerabilities with Code Property Graphs. \emph{IEEE S\&P}, 2014.
\bibitem{packer2023memgpt} C.~Packer et al. MemGPT: Towards LLMs as Operating
Systems. 2023. arXiv:2310.08560.
\bibitem{park2023generative} J.~S.~Park et al. Generative Agents: Interactive
Simulacra of Human Behavior. \emph{UIST}, 2023. arXiv:2304.03442.
\bibitem{langgraph} LangChain. LangGraph: Building Stateful, Multi-Actor
Applications with LLMs. Software framework, 2023--2024.
\url{https://github.com/langchain-ai/langgraph}.
\bibitem{denning1968working} P.~J.~Denning. The Working Set Model for Program
Behavior. \emph{Communications of the ACM}, 11(5), 1968.
\bibitem{liu2023lost} N.~F.~Liu et al. Lost in the Middle: How Language Models Use
Long Contexts. \emph{TACL}, 2024. arXiv:2307.03172.
\bibitem{kwon2023paged} W.~Kwon et al. Efficient Memory Management for Large
Language Model Serving with PagedAttention. \emph{SOSP}, 2023. arXiv:2309.06180.
\bibitem{haber1991timestamp} S.~Haber and W.~S.~Stornetta. How to Time-Stamp a
Digital Document. \emph{Journal of Cryptology}, 3(2), 1991.
\bibitem{jimenez2023swebench} C.~E.~Jimenez et al. SWE-bench: Can Language Models
Resolve Real-World GitHub Issues? \emph{ICLR}, 2024. arXiv:2310.06770.
\end{thebibliography}

\end{document}
